\documentclass[a4paper,14pt]{exam}
\usepackage{array,amsmath,fancybox,amsfonts,bm}
\usepackage{graphicx,amssymb,wasysym,latexsym,slashed,anysize,subfigure,calc}



%versores: %uso: \uveci\ne\hat{\imath}_{\uvec{i}}+\uvec{j}+\uvec{k}$
	\usepackage{bm}
	\newcommand{\uveci}{{\bm{\hat{\textnormal{\bfseries\i}}}}}
	\newcommand{\uvecj}{{\bm{\hat{\textnormal{\bfseries\j}}}}}
	\DeclareRobustCommand{\uvec}[1]{{%
 	 \ifcsname uvec#1\endcsname
  	   \csname uvec#1\endcsname
  	 \else
   		 \bm{\hat{\mathbf{#1}}}%
   		\fi
									}}



%tamaños fuentes extra: 8,9,14,17,20
	\usepackage{extsizes}
	
 %caja bonita
  \usepackage[many]{tcolorbox}
\definecolor{titlebg}{RGB}{80,100,72}
\newtcolorbox{learn}{
  breakable,
  enhanced,
  arc=0pt,
  outer arc=0pt,
  colframe=titlebg,
  colback=titlebg!7,
  overlay unbroken and first={
    \node[
      draw=titlebg,
      fill=titlebg,
      rotate=270,
      anchor=north west,
      text=white,
      font=\bfseries
    ]
    at (frame.north west)  
    {DATOS};
  }
}

%imágenes fundidas texto
\usepackage{wrapfig}

\usepackage{mwe,float,afterpage}

\usepackage{anysize,float}
	%\marginsize{left}{right}{top}{bottom}
    \marginsize{1.25cm}{1.25cm}{1cm}{1cm}
    
    
    
    
%entorno para química: texto y math mode
%\newcommand*\chem[1]{\textup{#1}}
\newcommand*\chem[1]{\ensuremath{\mathrm{#1}}}

%símbolo bonito para qed
\usepackage{amsthm}
	%\newcommand{\mc}[3]{\multicolumn{#1}{#2}{#3}}
\newcommand{\halmos}{\vspace*{-7mm}
	
					\begin{flushright}
						\rule{1mm}{2.5mm} 
					\end{flushright}\vspace*{-6mm}
						
					  }
\renewcommand{\qed}{\halmos}

\usepackage[spanish]{babel}
\usepackage[utf8]{inputenc}
%\usepackage{fancyhdr}
\usepackage{lettrine}

%subrayado y highlight. Tiene que estar detrás de usepackage[color]
\usepackage{soul}
	%\textst{normal} 
	\DeclareRobustCommand{\hlcyan}[1]{{\sethlcolor{cyan}\hl{#1}}} %en cyan!
		%\hl{highlighted text} %uso


%fuentes
	\usepackage[T1]{fontenc}
	\usepackage{mathpazo}

	\usepackage{librebaskerville}

%fancy tables
	\usepackage{booktabs}
	\usepackage[small,centerlast,up]{caption}

	

\usepackage{amssymb,amsmath}
\newcommand{\angstrom}{\mbox{\normalfont\AA}}

\usepackage{graphicx,graphics}
\usepackage{color}
	%	\colorgrids %grids con color
	%	\definecolor{GridColor}{gray}{0.8}
		%grid
%			\setlength{\gridsize}{5mm}
%   			 \setlength{\gridlinewidth}{0.1pt}

%fracciones en texto más elaboradas
	\usepackage{xfrac}

%mejora en la exposición de guarismos
	\usepackage{numprint}
 


%para los entornos de soluciones
\usepackage{mdframed,xcolor,color}
	%uso: 	entorno (begin,end): {mdframed}[backgroundcolor=brown!10] 


\usepackage{everypage,draftwatermark} %Marca de agua
\SetWatermarkText{\parbox{2\textwidth}{\centering B2\\dist.}}
\SetWatermarkScale{1.5}
\SetWatermarkLightness{0.92}





\pagestyle{headandfoot}
%\runningheadrule
%\firstpageheader{Física para Biólogos}{Boletín 1a}{1 de octubre de 2011}
%\runningheader{Física para Biólogos}
%              {Boletín 1a, Página \thepage\ de \numpages}
%            {1 de octubre de 2011}
%\firstpagefooter{}{}{}
%\runningfooter{}{}{}


%espacios extra
%\extrawidth{1in}
%%\extrafootheight{.75in}
%\extraheadheight{1in}

%cabecera
\header{\includegraphics[scale=.22]{a.png}}       
      {}
       {$f_2$\textbf{B2-{\small distancia}}\\\textbf{febrero}{\small 2023}}       
%pie de página

   %enumera en todas menos en la primera
   %\runningfootrule

   %enumera en todas
   %\footrule
   %\footer{}{Página \thepage\ de \numpages}{}
   \footer{}
	{%Página \thepage\ de \numpages
	}
	{%\iflastpage{{\footnotesize \textit{Fin del examen}}}{{\footnotesize \textit{Da la vuelta a la hoja}\ldots}}
	}

	
\begin{document}
	\pagenumbering{gobble}
	
\vspace*{-2ex}

\fbox{\parbox{0.945\linewidth}{\textbf{Nombre y apellidos}:\vspace*{2mm}}}
\vspace{3ex}

\begin{enumerate}
	
	
		\item   Un objeto de $5\;cm$ de altura se encuentra a una distancia $s$ de una lente delgada 
convergente. La lente forma una imagen real e invertida del objeto. El tamaño de la imagen es de
$10\;cm$. La distancia focal de la lente es $10\;cm$.
			\begin{itemize}
				\item[a)] Determine la distancia a la cual se encuentra el objeto de la lente.
				\item[b)] Realice el diagrama de rayos del sistema.				
			\end{itemize}
			\qed
			
			
	\item Determina la posición donde debe colocarse un objeto real situado a la izquierda de una lente delgada convergente de $3$ dioptrías para que el tamaño de la imagen formada  sea:
			\begin{itemize}
				\item[a)] Recta y de tamaño doble que el objeto.
				\item[b)] Invertida y de mitad de tamaño que el objeto.
			\end{itemize}
%				\begin{small}
%					\textbf{Dato:} 
%				\end{small}
		\qed
	
	

	
	
	
	\item 	 Un sistema óptico está formado por dos lentes convergentes delgadas de
distancia focal $50\;mm$, cada una, separadas una distancia de $200\;mm$. Un objeto de $20\;mm$ de
altura está situado a una distancia de $100\;mm$ delante de la primera lente.
			\begin{itemize}
				\item[a)] Realice el trazado de rayos de la imagen formada por la primera lente y calcule su posición.
				\item[b)] Realice el trazado de rayos de la imagen formada por el sistema óptico completo y calcule su posición.
			\end{itemize}
	\qed
	
	
	

			\item Tres cargas gravitatorias iguales y de valor $m$ yacen, estáticas, en sendos vértices de un cuadrado de lado $\ell$, tal y como se indica en la figura adjunta. 
	
 		\begin{minipage}{.7\linewidth}
 		\begin{itemize}
			\item[a)] Calcula el campo gravitatorio provocado en el punto $\mathcal{P}$ correspondiente al centro del cuadrado.
			\item[b)] Calcula la energía potencial gravitatoria que tendría una masa $M$ que se colocara en el punto $\mathcal{Q}$, correspondiente al vértice libre.
			\vspace{.1ex}
						
		\end{itemize}
 		\end{minipage}\hfill
 		\begin{minipage}{.3\linewidth}
 		\raggedleft
 			\def\svgwidth{.85\linewidth}
			\input{triangulo.pdf_tex}
 		\end{minipage}
 		\begin{small}
				
			\textbf{Datos:} $G\simeq 6.67\cdot 10^{-11}\;Nm^2kg^{-2},\hspace{3mm}m=1\;kg,\hspace{3mm}M=2\,kg,\hspace{3mm}\ell=1\,m$.
			\end{small}
			
		\qed
	
		
\end{enumerate}		
 
\end{document}

